\documentclass[12pt,tikz,border=5pt]{article}
\usepackage[utf8]{inputenc}
\usepackage[spanish]{babel}
\addto\captionsspanish{\renewcommand{\tablename}{Tabla}}
\decimalpoint
\usepackage{amsmath, amssymb, amsthm, amsfonts}
\usepackage{geometry}
\usepackage{enumitem}
\usepackage{multicol}
\usepackage{tikz}
\usepackage{subcaption}
\usetikzlibrary{babel} 
\usepackage{setspace}
\usepackage{booktabs}
\usepackage{natbib}   
\usepackage{hyperref}
\geometry{margin=2.5cm}
\newtheorem{ejercicio}{Ejercicio}[section]

\title{Cálculo Diferencial e Integral \\ \large Licenciatura en Ingeniería en Computación}
\author{M. en C. Edwin Pérez \\ Dr. Julian Bravo}
\date{20 de octubre de 2026}

\begin{document}
\maketitle

\section{Derivabilidad y continuidad}

\noindent \textbf{¿Existe alguna relación entre derivabilidad y continuidad?}

Recordemos que una manera de expresar la derivada de $f$ en un punto $a$ es

$$f'(a) = \lim_{x\to a} \frac{f(x)-f(a)}{x-a}$$

\textsc{siempre que dicho límite exista.}

La existencia de dicho límite requiere conocer las siguientes definiciones:

\subsection*{Definición: Derivada Lateral}

\begin{enumerate}[label=\alph*)]
\item Si la función $f$ está definida en un punto $a$, entonces la
\textbf{derivada por la derecha} de $f$ en $a$, denotada por $f'_+(a)$, está
definida por
$$f'_+(a) = \lim_{x\to a^+} \frac{f(x)-f(a)}{x-a}$$

\item Si la función $f$ está definida en un punto $a$, entonces la
\textbf{derivada por la izquierda} de $f$ en $a$, denotada por $f'_-(a)$,
está definida por
$$f'_-(a) = \lim_{x\to a^-} \frac{f(x)-f(a)}{x-a}$$
\end{enumerate}

Entonces, podemos concluir lo siguiente:

\begin{quote}
El $\displaystyle\lim_{x\to a}\frac{f(x)-f(a)}{x-a} = f'(a)$ existe si y
sólo si $f'_+(a)$, $f'_-(a)$ existen y $f'_+(a) = f'_-(a)$.
\end{quote}

Conocer lo anterior, nos permite definir lo siguiente:

\subsection*{Definición}

Una función $f$ es \textbf{derivable en $a$} si $f'(a)$ existe. Es
\textbf{derivable en un intervalo abierto $(a,b)$}\footnote{También puede
ser $(a,\infty)$, $(-\infty,a)$ o $(-\infty,\infty)$.} si es derivable
$\forall\, c\in(a,b)$.

\section{Ejemplos ilustrativos}

\begin{ejercicio}
¿Dónde es derivable la función $f(x) = \dfrac{3}{x}$?

Primero calculemos la derivada de $f$:

\begin{align*}
f'(x) &= \lim_{\Delta x\to 0}\frac{f(x+\Delta x)-f(x)}{\Delta x} = \lim_{\Delta x\to 0}\frac{\dfrac{3}{x+\Delta x}-\dfrac{3}{x}}{\Delta x} \\[4pt]
&= \lim_{\Delta x\to 0}\frac{\dfrac{3x-3(x+\Delta x)}{x(x+\Delta x)}}{\Delta x} = \lim_{\Delta x\to 0}\frac{3x-3x-3\Delta x}{x^2\Delta x+x\Delta x^2} \\[4pt]
&= \lim_{\Delta x\to 0}\frac{-3\Delta x}{\Delta x(x^2+x\Delta x)} = \lim_{\Delta x\to 0}\frac{-3}{x^2+x\Delta x} = -\frac{3}{x^2}
\end{align*}

Notemos que $\operatorname{Dom} f' = (-\infty,0)\cup(0,\infty)$, por lo que
$f$ es derivable $\forall x\in(-\infty,0)\cup(0,\infty)$.

Además observemos que $\operatorname{Dom} f = \operatorname{Dom} f'$.
\end{ejercicio}

\begin{ejercicio}
¿La función $g(x) = \sqrt{x}$ es derivable en $x=0$?

Para valores de $x<0$, $g$ no está definida; por lo tanto, el cociente $\frac{g(0+\Delta x)-g(0)}{\Delta x}$ no tiene sentido cuando $\Delta x<0$, y $g'_-(0)$ no existe.
Ahora veamos qué sucede con la derivada lateral derecha de $g$ en 0:

\begin{align*}
g'_+(0) &= \lim_{\Delta x\to 0^+}\frac{f(0+\Delta x)-f(0)}{\Delta x} = \lim_{\Delta x\to 0^+}\frac{\sqrt{\Delta x}-\sqrt{0}}{\Delta x} \\[4pt]
&= \lim_{\Delta x\to 0^+}\frac{\sqrt{\Delta x}}{\Delta x} = \lim_{\Delta x\to 0^+}\frac{1}{\sqrt{\Delta x}} \qquad \text{¡No existe!}
\end{align*}

Así que $g'_{+}(0)$ no existe. \textbf{Conclusión}: $g$ no es derivable en $x=0$. \bigskip

¿En qué intervalo la función $g$ es derivable? Primero calculemos la derivada de $g$:

\begin{align*}
g'(x) &= \lim_{\Delta x\to 0}\frac{f(x+\Delta x)-f(x)}{\Delta x} = \lim_{\Delta x\to 0}\frac{\sqrt{x+\Delta x}-\sqrt{x}}{\Delta x} \\[4pt]
&= \lim_{\Delta x\to 0}\frac{\sqrt{x+\Delta x}-\sqrt{x}}{\Delta x}\cdot\left(\frac{\sqrt{x+\Delta x}+\sqrt{x}}{\sqrt{x+\Delta x}+\sqrt{x}}\right) \\[4pt]
&= \lim_{\Delta x\to 0}\frac{(x+\Delta x)-x}{\Delta x\sqrt{x+\Delta x}+\Delta x\sqrt{x}} = \lim_{\Delta x\to 0}\frac{\Delta x}{\Delta x\sqrt{x+\Delta x}+\Delta x\sqrt{x}} \\[4pt]
&= \lim_{\Delta x\to 0}\frac{1}{\sqrt{x+\Delta x}+\sqrt{x}} = \frac{1}{2\sqrt{x}}
\end{align*}

 De modo que el $\operatorname{Dom} g' = (0,\infty)$. Por lo que $g$ es derivable $\forall x\in(0,\infty)$. Sabemos que el $\operatorname{Dom} g = [0,\infty)$. Luego, $\operatorname{Dom} g' \subset \operatorname{Dom} g$.

\end{ejercicio}

\subsection*{iii) $h(x) = |x|$}

Recordemos que la función $h$ se puede reescribir como

$$h(x) = \begin{cases} x & \text{si } x\ge 0 \\ -x & \text{si } x<0 \end{cases}$$

Por un lado, considere los valores de $x>0$. Si se considera un $\Delta x$
lo suficientemente pequeño para que $x+\Delta x>0$, entonces
$|x+\Delta x| = x+\Delta x$. Así que

$$h'(x) = \lim_{\Delta x\to 0}\frac{h(x+\Delta x)-h(x)}{\Delta x} = \lim_{\Delta x\to 0}\frac{|x+\Delta x|-|x|}{\Delta x} = \lim_{\Delta x\to 0}\frac{x+\Delta x-x}{\Delta x} = \lim_{\Delta x\to 0}\frac{\Delta x}{\Delta x} = 1$$

Por lo tanto, $h$ es derivable $\forall x\in(0,\infty)$.

Por otro lado, considere aquellos valores para los cuales $x<0$. Elijamos
$\Delta x$ lo suficientemente pequeño para que $x+\Delta x<0$, entonces
$|x+\Delta x| = -(x+\Delta x)$. Entonces,

$$h'(x) = \lim_{\Delta x\to 0}\frac{|x+\Delta x|-|x|}{\Delta x} = \lim_{\Delta x\to 0}\frac{-(x+\Delta x)+x}{\Delta x} = -\lim_{\Delta x\to 0}\frac{\Delta x}{\Delta x} = -1$$

Por lo tanto, $h$ es derivable $\forall x\in(-\infty,0)$.

Por último, considere el caso cuando $x=0$. Observemos que

$$h'(0) = \lim_{\Delta x\to 0}\frac{h(0+\Delta x)-h(0)}{\Delta x} = \lim_{\Delta x\to 0}\frac{|\Delta x|}{\Delta x} \qquad \text{¿Existirá?}$$

Utilicemos la derivada lateral para descubrir si el límite anterior existe:
por un lado,

$$h'_+(0) = \lim_{\Delta x\to 0^+}\frac{|\Delta x|}{\Delta x} = \lim_{\Delta x\to 0^+} 1 = 1$$

por otro lado,

$$h'_-(0) = \lim_{\Delta x\to 0^-}\frac{|\Delta x|}{\Delta x} = \lim_{\Delta x\to 0^-}(-1) = -1$$

Debido a que $h'_-(0) \ne h'_+(0)$, $h$ no es derivable en $0$.

\textbf{En conclusión}, $h$ es derivable $\forall x\in(-\infty,0)\cup(0,\infty)$.

\subsection*{iv) $f(x) = \sqrt[3]{x}$}

\noindent\textbf{¿Por qué $f$ no es derivable en $x=0$?} Observemos que

$$f'(0) = \lim_{\Delta x\to 0}\frac{f(0+\Delta x)-f(0)}{\Delta x} = \lim_{\Delta x\to 0}\frac{\Delta x^{1/3}}{\Delta x} = \lim_{\Delta x\to 0}\frac{1}{\Delta x^{2/3}} \qquad \text{¡No existe!}$$

Por lo tanto, $f$ no es derivable en $0$.

\bigskip

Tanto la continuidad como la derivabilidad son propiedades deseables para
una función, y el siguiente teorema muestra cómo se relacionan ambas.

\section{Relación entre derivabilidad y continuidad}

\textbf{Teorema.} Si $f$ es derivable en $a$, entonces $f$ es continua en
$a$.

\textbf{Demostración.}

Recordemos que una función $f$ es continua en $a$ si y sólo si cumple con
lo siguiente:

\begin{enumerate}[label=\roman*)]
\item $f(a)$ existe
\item $\displaystyle\lim_{x\to a} f(x)$ existe
\item $\displaystyle\lim_{x\to a} f(x) = f(a)$
\end{enumerate}

Como $f$ es derivable en $a$, esto quiere decir que

$$f'(a) = \lim_{x\to a}\frac{f(x)-f(a)}{x-a} \quad \text{existe}$$

por lo que $f(a)$ existe y se cumple con el i) de las condiciones de
continuidad.

Para demostrar que $\displaystyle\lim_{x\to a} f(x)$ existe y
$\displaystyle\lim_{x\to a} f(x) = f(a)$, basta con demostrar que

$$\lim_{x\to a}\big[f(x)-f(a)\big] = 0 \qquad \text{¿Por qué?}$$

ya que $\displaystyle\lim_{x\to a} f(x) = f(a)$ si y sólo si
$\displaystyle\lim_{x\to a}\big[f(x)-f(a)\big] = 0$. Entonces

\begin{align*}
\lim_{x\to a}\big[f(x)-f(a)\big] &= \lim_{x\to a}\big[\big(f(x)-f(a)\big)\cdot 1\big] = \lim_{x\to a}\left[\frac{f(x)-f(a)}{x-a}\cdot(x-a)\right] \\[4pt]
&= \lim_{x\to a}\frac{f(x)-f(a)}{x-a}\cdot \lim_{x\to a}(x-a) \\[4pt]
&= f'(a)\cdot 0 = 0 \qquad \blacksquare
\end{align*}

\section{Reglas de Derivación}


\subsection{Regla de la suma (y de la resta)}

\textbf{Teorema.} Si $f$ y $g$ son derivables en $x$, entonces

$$(f+g)'(x) = f'(x) + g'(x) \qquad \text{y} \qquad (f-g)'(x) = f'(x) - g'(x)$$

\textbf{Demostración.}

Por definición de derivada, aplicada a la función $f+g$:

\begin{align*}
(f+g)'(x) &= \lim_{\Delta x\to 0} \frac{(f+g)(x+\Delta x)-(f+g)(x)}{\Delta x} \\[4pt]
&= \lim_{\Delta x\to 0} \frac{\big[f(x+\Delta x)+g(x+\Delta x)\big]-\big[f(x)+g(x)\big]}{\Delta x}
\end{align*}

Reagrupamos los términos de $f$ y de $g$ por separado:

$$(f+g)'(x) = \lim_{\Delta x\to 0}\left[\frac{f(x+\Delta x)-f(x)}{\Delta x} + \frac{g(x+\Delta x)-g(x)}{\Delta x}\right]$$

Como $f$ y $g$ son derivables, ambos límites existen por separado, así que
podemos separar el límite de la suma en la suma de los límites:

$$(f+g)'(x) = \lim_{\Delta x\to 0}\frac{f(x+\Delta x)-f(x)}{\Delta x} + \lim_{\Delta x\to 0}\frac{g(x+\Delta x)-g(x)}{\Delta x} = f'(x)+g'(x) \qquad \blacksquare$$

La demostración de la resta es idéntica, cambiando cada signo $+$ por un
signo $-$.

\subsection{Regla del producto}

\textbf{Teorema.} Si $f$ y $g$ son derivables en $x$, entonces

$$(fg)'(x) = f'(x)\,g(x) + f(x)\,g'(x)$$

\textbf{Demostración.}

Por definición:

$$(fg)'(x) = \lim_{\Delta x\to 0}\frac{f(x+\Delta x)g(x+\Delta x) - f(x)g(x)}{\Delta x}$$

Notemos que no podemos separar directamente este límite como hicimos con la
suma. El truco aquí es \textbf{sumar y restar} el mismo término,
$f(x+\Delta x)\,g(x)$, en el numerador — esto no cambia el valor de la
expresión, pero nos permite factorizar:

\begin{align*}
(fg)'(x) &= \lim_{\Delta x\to 0}\frac{f(x+\Delta x)g(x+\Delta x) - f(x+\Delta x)g(x) + f(x+\Delta x)g(x) - f(x)g(x)}{\Delta x} \\[4pt]
&= \lim_{\Delta x\to 0}\left[f(x+\Delta x)\cdot\frac{g(x+\Delta x)-g(x)}{\Delta x} + g(x)\cdot\frac{f(x+\Delta x)-f(x)}{\Delta x}\right]
\end{align*}

Como $f$ es derivable, por el teorema de la sección anterior sabemos que
$f$ también es continua, así que $f(x+\Delta x)\to f(x)$ cuando
$\Delta x\to 0$. Tomando el límite de cada factor por separado:

$$(fg)'(x) = f(x)\,g'(x) + g(x)\,f'(x) \qquad \blacksquare$$

\subsection{Regla del cociente}

\textbf{Teorema.} Si $u$ y $v$ son derivables en $x$, con $v(x)\ne 0$,
entonces

$$\left(\frac{u}{v}\right)'(x) = \frac{u'(x)\,v(x) - u(x)\,v'(x)}{\big(v(x)\big)^2}$$

\textbf{Demostración.}

Por definición:

$$\left(\frac{u}{v}\right)'(x) = \lim_{\Delta x\to 0}\frac{1}{\Delta x}\left[\frac{u(x+\Delta x)}{v(x+\Delta x)} - \frac{u(x)}{v(x)}\right]$$

Juntamos las dos fracciones de en medio con un común denominador:

$$= \lim_{\Delta x\to 0}\frac{u(x+\Delta x)v(x) - u(x)v(x+\Delta x)}{\Delta x\,v(x+\Delta x)v(x)}$$

Igual que con el producto, sumamos y restamos un término conveniente en el
numerador, esta vez $u(x)v(x)$:

$$= \lim_{\Delta x\to 0}\frac{\big[u(x+\Delta x)-u(x)\big]v(x) - u(x)\big[v(x+\Delta x)-v(x)\big]}{\Delta x\,v(x+\Delta x)v(x)}$$

Separamos en dos fracciones y dividimos cada numerador entre $\Delta x$:

$$= \lim_{\Delta x\to 0}\frac{\dfrac{u(x+\Delta x)-u(x)}{\Delta x}\,v(x) - u(x)\,\dfrac{v(x+\Delta x)-v(x)}{\Delta x}}{v(x+\Delta x)\,v(x)}$$

Como $v$ es derivable, también es continua, así que $v(x+\Delta x)\to v(x)$
cuando $\Delta x\to 0$. Tomando el límite de cada parte:

$$\left(\frac{u}{v}\right)'(x) = \frac{u'(x)v(x) - u(x)v'(x)}{\big(v(x)\big)^2} \qquad \blacksquare$$

\subsection{Regla de la cadena}

\textbf{Teorema.} Sea $y=f(u)$ y $u=g(x)$, con $g$ derivable en $x$ y $f$
derivable en $g(x)$. Entonces

$$\frac{dy}{dx} = \frac{dy}{du}\cdot\frac{du}{dx}$$

\textbf{Demostración (heurística).}

Sea $u=g(x)$, y sea $\Delta u = g(x+\Delta x)-g(x)$ el cambio
correspondiente en $u$ cuando $x$ cambia en $\Delta x$. Si $\Delta u\ne 0$,
podemos multiplicar y dividir por $\Delta u$ dentro del cociente de
diferencias de $y=f(u)=f(g(x))$:

$$\frac{\Delta y}{\Delta x} = \frac{\Delta y}{\Delta u}\cdot\frac{\Delta u}{\Delta x}$$

Cuando $\Delta x\to 0$, también $\Delta u\to 0$ (porque $g$ es derivable, y
por lo tanto continua). Entonces, usando la propiedad del límite de un
producto:

$$\frac{dy}{dx} = \lim_{\Delta x\to 0}\frac{\Delta y}{\Delta x} = \left(\lim_{\Delta u\to 0}\frac{\Delta y}{\Delta u}\right)\cdot\left(\lim_{\Delta x\to 0}\frac{\Delta u}{\Delta x}\right) = \frac{dy}{du}\cdot\frac{du}{dx} \qquad \blacksquare$$

\textbf{Advertencia técnica.} Este argumento tiene un hueco: $\Delta u$
podría valer $0$ para algunos $\Delta x\ne 0$ arbitrariamente pequeños (por
ejemplo, si $g$ es constante en ciertos tramos), y entonces no se podría
dividir entre $\Delta u$. La demostración rigurosa, que evita este
problema, requiere definir con más cuidado el cociente usando una función
auxiliar continua en $0$; queda fuera del alcance de este curso. El
resultado, sin embargo, sí es válido en general.

\subsection{Composición de funciones}

\textbf{Definición.} Dadas dos funciones $f$ y $g$, su \textbf{composición}
$f\circ g$ se define por

$$(f\circ g)(x) = f\big(g(x)\big)$$

Es decir, primero se evalúa $g$ en $x$, y después se evalúa $f$ en el
resultado obtenido. El dominio de $f\circ g$ consta de los valores de $x$
en el dominio de $g$ tales que $g(x)$ esté en el dominio de $f$.

\medskip

La regla de la cadena de la sección anterior, escrita en términos de la
composición $f\circ g$, toma la forma:

$$(f\circ g)'(x) = f'\big(g(x)\big)\cdot g'(x)$$

Es decir: la derivada de la composición es la derivada de la función
``de afuera'' ($f$), evaluada en la función ``de adentro'' ($g(x)$),
multiplicada por la derivada de la función ``de adentro''.

\medskip

\textbf{Ejemplo ilustrativo.} Sea $F(x) = (f\circ g)(x)$, con
$f(u)=u^2$ y $g(x)=3x+1$, de modo que $F(x) = (3x+1)^2$.

Usando la regla de la cadena: $f'(u)=2u$, así que $f'(g(x)) = 2(3x+1)$; y
$g'(x)=3$. Entonces,

$$F'(x) = f'\big(g(x)\big)\cdot g'(x) = 2(3x+1)\cdot 3 = 6(3x+1) = 18x+6$$

\textbf{Verificación:} expandiendo directamente, $F(x)=(3x+1)^2 = 9x^2+6x+1$,
y derivando término por término (regla de la potencia y de la suma):
$F'(x) = 18x+6$. Coincide con el resultado obtenido por la regla de la
cadena. $\blacksquare$

\end{document}

\end{document}
